\documentclass[10pt,a4paper]{article}

%double si on veut une version prof et une version élève. simple si on veut une seule version.
\newcommand{\buildmode}{simple}

\InputIfFileExists{version.tex}{}{}
% Version (définie par le script)
\providecommand{\version}{eleve}

\usepackage{feuilleinterro}
\usepackage{schemaevolution}

%%%à modifier à chaque fois

\renewcommand{\niveau}{1M}
\renewcommand{\chapitre}{Évolutions}
\renewcommand{\numchap}{0.3}
\renewcommand{\theme}{Statistiques}

%%%titre "CORRECTION" au lieu de "INTERROGATION", sans les champs Nom/Prénom
\renewcommand{\interroversion}[1]{
    \noindent\textbf{\niveau}

    \vspace{-0.5cm}
    \begin{center}
        \textbf{\large CORRECTION -- \chapitre}
    \end{center}

    \vspace{0.1cm}
}

\begin{document}

% =========================================================
% PAGE 1 : VERSION A
% =========================================================

\interroversion{A}

\textbf{Exercice 1 -- Solution d'une équation \hfill /2}

Le nombre \(3\) est-il solution de l'équation \(2x+3=5x-6\) ?

On remplace \(x\) par \(3\) dans chaque membre de l'équation :
\begin{itemize}
\item membre de gauche : \(2\times3+3=9\) ;
\item membre de droite : \(5\times3-6=9\).
\end{itemize}
Les deux membres sont égaux, donc \textcolor{blue}{\(3\) est solution} de l'équation.

\textbf{Exercice 2 -- Résolution d'une équation \hfill /3}

\[
\begin{aligned}
6x+5&=9\\
6x&=9-5\\
6x&=4\\
x&=\frac{4}{6}\\
x&=\frac{2}{3}.
\end{aligned}
\]

La solution de l'équation est \(\textcolor{blue}{\dfrac{2}{3}}\).

\textbf{Exercice 3 -- Valeur finale \hfill /3}

Le coefficient multiplicateur associé à une évolution de \(-50\%\) est \(1-\dfrac{50}{100}=0{,}5\).

Valeur finale : \(60\times0{,}5=\textcolor{blue}{30}\).

\schemaevolution{60}{\textcolor{blue}{30}}{-50\%}{\textcolor{blue}{\times 0{,}5}}

\textbf{Exercice 4 -- Taux d'évolution \hfill /2}

\[
t=\frac{V_F-V_I}{V_I}=\frac{100-80}{80}=\frac{20}{80}=\frac{1}{4}=0{,}25=\textcolor{blue}{+25\%}.
\]

Il s'agit d'une augmentation de 25\%.

\schemaevolution{80}{100}{\textcolor{blue}{+25\%}}{}

\newpage

% =========================================================
% PAGE 2 : VERSION B
% =========================================================

\interroversion{B}

\textbf{Exercice 1 -- Solution d'une équation \hfill /2}

Le nombre \(4\) est-il solution de l'équation \(3x-2=x+5\) ?

On remplace \(x\) par \(4\) dans chaque membre de l'équation :
\begin{itemize}
\item membre de gauche : \(3\times4-2=10\) ;
\item membre de droite : \(4+5=9\).
\end{itemize}
Les deux membres ne sont pas égaux (\(10\neq9\)), donc \textcolor{blue}{\(4\) n'est pas solution} de l'équation.

\textbf{Exercice 2 -- Résolution d'une équation \hfill /3}

\[
\begin{aligned}
4x-3&=7\\
4x&=7+3\\
4x&=10\\
x&=\frac{10}{4}\\
x&=\frac{5}{2}.
\end{aligned}
\]

La solution de l'équation est \(\textcolor{blue}{\dfrac{5}{2}}\).

\textbf{Exercice 3 -- Valeur finale \hfill /3}

Le coefficient multiplicateur associé à une évolution de \(+50\%\) est \(1+\dfrac{50}{100}=1{,}5\).

Valeur finale : \(40\times1{,}5=\textcolor{blue}{60}\).

\schemaevolution{40}{\textcolor{blue}{60}}{+50\%}{\textcolor{blue}{\times 1{,}5}}

\textbf{Exercice 4 -- Taux d'évolution \hfill /2}

\[
t=\frac{V_F-V_I}{V_I}=\frac{40-50}{50}=-\frac{10}{50}=-\frac{1}{5}=-0{,}2=\textcolor{blue}{-20\%}.
\]

Il s'agit d'une diminution de 20\%.

\schemaevolution{50}{40}{\textcolor{blue}{-20\%}}{}

\newpage

% =========================================================
% PAGE 3 : VERSION C
% =========================================================

\interroversion{C}

\textbf{Exercice 1 -- Solution d'une équation \hfill /2}

Le nombre \(5\) est-il solution de l'équation \(4x-7=2x+3\) ?

On remplace \(x\) par \(5\) dans chaque membre de l'équation :
\begin{itemize}
\item membre de gauche : \(4\times5-7=13\) ;
\item membre de droite : \(2\times5+3=13\).
\end{itemize}
Les deux membres sont égaux, donc \textcolor{blue}{\(5\) est solution} de l'équation.

\textbf{Exercice 2 -- Résolution d'une équation \hfill /3}

\[
\begin{aligned}
10x+1&=5\\
10x&=5-1\\
10x&=4\\
x&=\frac{4}{10}\\
x&=\frac{2}{5}.
\end{aligned}
\]

La solution de l'équation est \(\textcolor{blue}{\dfrac{2}{5}}\).

\textbf{Exercice 3 -- Valeur finale \hfill /3}

Le coefficient multiplicateur associé à une évolution de \(-50\%\) est \(1-\dfrac{50}{100}=0{,}5\).

Valeur finale : \(80\times0{,}5=\textcolor{blue}{40}\).

\schemaevolution{80}{\textcolor{blue}{40}}{-50\%}{\textcolor{blue}{\times 0{,}5}}

\textbf{Exercice 4 -- Taux d'évolution \hfill /2}

\[
t=\frac{V_F-V_I}{V_I}=\frac{250-200}{200}=\frac{50}{200}=\frac{1}{4}=0{,}25=\textcolor{blue}{+25\%}.
\]

Il s'agit d'une augmentation de 25\%.

\schemaevolution{200}{250}{\textcolor{blue}{+25\%}}{}

\newpage

% =========================================================
% PAGE 4 : VERSION D
% =========================================================

\interroversion{D}

\textbf{Exercice 1 -- Solution d'une équation \hfill /2}

Le nombre \(2\) est-il solution de l'équation \(5x+1=3x+4\) ?

On remplace \(x\) par \(2\) dans chaque membre de l'équation :
\begin{itemize}
\item membre de gauche : \(5\times2+1=11\) ;
\item membre de droite : \(3\times2+4=10\).
\end{itemize}
Les deux membres ne sont pas égaux (\(11\neq10\)), donc \textcolor{blue}{\(2\) n'est pas solution} de l'équation.

\textbf{Exercice 2 -- Résolution d'une équation \hfill /3}

\[
\begin{aligned}
8x-5&=1\\
8x&=1+5\\
8x&=6\\
x&=\frac{6}{8}\\
x&=\frac{3}{4}.
\end{aligned}
\]

La solution de l'équation est \(\textcolor{blue}{\dfrac{3}{4}}\).

\textbf{Exercice 3 -- Valeur finale \hfill /3}

Le coefficient multiplicateur associé à une évolution de \(+50\%\) est \(1+\dfrac{50}{100}=1{,}5\).

Valeur finale : \(50\times1{,}5=\textcolor{blue}{75}\).

\schemaevolution{50}{\textcolor{blue}{75}}{+50\%}{\textcolor{blue}{\times 1{,}5}}

\textbf{Exercice 4 -- Taux d'évolution \hfill /2}

\[
t=\frac{V_F-V_I}{V_I}=\frac{15-20}{20}=-\frac{5}{20}=-\frac{1}{4}=-0{,}25=\textcolor{blue}{-25\%}.
\]

Il s'agit d'une diminution de 25\%.

\schemaevolution{20}{15}{\textcolor{blue}{-25\%}}{}

\end{document}
